\subsection{Proveedor de Tiempo}

\subsubsection{Regla: el tiempo se obtiene de un \texttt{TimeProvider} inyectado}

\textbf{Regla:} Prohibido leer el reloj del sistema con \texttt{DateTime.Now}, \texttt{DateTime.UtcNow} o \texttt{DateTimeOffset.UtcNow} fuera de la raíz de composición. Una clase que necesita la hora actual recibe un \texttt{TimeProvider} por constructor y llama a \texttt{GetUtcNow()}. En la raíz de composición se registra \texttt{TimeProvider.System}.

\textbf{Descripción:} El tiempo es un resultado no determinista y, por lo tanto, un servicio externo (Sección 10.1). \texttt{TimeProvider} es la abstracción del tiempo que ofrece .NET, y su guía lo justifica así:
\begin{quote}
By using \texttt{TimeProvider}, you ensure that your code is testable and predictable. (Microsoft, 2026j, párr. 1)
\end{quote}
\noindent Para las pruebas, el paquete \texttt{Microsoft.\allowbreak Extensions.\allowbreak TimeProvider.\allowbreak Testing} ofrece \texttt{FakeTimeProvider}, una implementación que se puede controlar (Microsoft, 2026j, sección FakeTimeProvider implementation, párr. 1).

\textbf{Código correcto:}
\begin{lstlisting}[style=csharpstyle]
public sealed class TurnTimer
{
    private readonly TimeProvider _timeProvider;

    public TurnTimer(TimeProvider timeProvider)
    {
        _timeProvider = timeProvider;
    }

    public bool HasTurnExpired(DateTimeOffset turnDeadline)
    {
        return _timeProvider.GetUtcNow() > turnDeadline;
    }
}
\end{lstlisting}

\textbf{Código incorrecto:}
\begin{lstlisting}[style=csharpstyle]
public sealed class TurnTimer
{
    public bool HasTurnExpired(DateTimeOffset turnDeadline)
    {
        return DateTimeOffset.UtcNow > turnDeadline;
    }
}
\end{lstlisting}
